\section{Synthesis}
\label{sec:synthesis}

Read together, the five works tell a story that none of them tells alone.
The statistic barely changed: from Peirce's residuals to Grubbs' studentised
deviate to the modified z-score to a forecast innovation, the measure of
surprise has been a scaled distance from expectation for 170 years. What
changed, era by era, is the expectation. Theory supplied it first, in the
form of a parametric null; robust statistics let the data supply their own
resistant version of it; the operational era replaced distributions with
engineering knowledge, instrument limits, climatology and the testimony of
neighbouring sensors; the framework era made the expectation institutional,
a versioned archive of what healthy instruments have done before; and the
current era learns it, from models and labelled history. Automated quality
control matured not by inventing better mathematics of surprise but by
enriching the definition of the unsurprising.

The second thread is less flattering and more useful. At every stage, the
threshold, the constant that converts a score into a verdict, has been the
component with the weakest justification and the longest life. Grubbs'
alpha was principled but assumed a world without autocorrelation; the 3.5 of
Iglewicz and Hoaglin was calibrated for batches and then wandered into
streams; the Mesonet's constants were judgements that hardened into
defaults; and my own experiments show, on a benchmark where the truth is
known, that no threshold can rescue a statistic from a broken assumption,
while window, threshold, sampling interval and noise scale must be designed
together or not at all. If I could add one sentence to every paper on this
list, it would be the same sentence: report what your thresholds cost, in
false alarms per reviewer per day and in faults foregone, because that
premium, Anscombe's premium, is the real currency of monitoring.

The third thread is architectural, and it is the lineage's quiet triumph.
Everything that survived is a composition: batteries rather than single
tests, because Figure~\ref{fig:detectors} is what single tests look like;
flags rather than deletions, because doubt is information; graded data
products rather than one truth, because users differ in their tolerance for
being wrong; and a human at the end, not because algorithms are weak but
because a flag is a message, and messages need readers. The 2019 framework,
for all its statistical sophistication, ships the same architecture as the
2000 one. I expect the learning era to do the same, and the works on this
list are why I expect it.

What I take back to my own work is concrete. The variance form of the
persistence test, in the reading list since 2000, dominates the equality
form my generation reaches for by reflex. The spatial check is the only
member of the battery that sees slow drift, and its latency is the price of
that sight, so it must be budgeted, not resented. The modified z-score is a
fine statistic attached to a constant that must be re-derived for every
sampling regime, and I now know what re-deriving it looks like. And every
constant I set from now on gets a written provenance, because I have spent a
semester on the other side of that omission, reverse-engineering the
judgements of people who, like the authors on this list, were sure at the
time that their reasons were too obvious to record.

Deciding what to doubt turns out to be the least automatable step in an
automated system. The five works on this reading list built, over 170
years, the machinery that makes the decision cheap, repeatable and mostly
right; what they hand to my generation is the remaining, harder task of
making it accountable.
